\documentclass[12pt]{article}
\usepackage[utf8]{inputenc}
\usepackage[T1]{fontenc}
\usepackage{hyperref}
\title{Buyang, E'ma Language Resources for K-12 Curriculum Development}
\author{Compiled by Claude Sonnet 5.5 (claude-sonnet-5-5) for Andr\'as Kornai}
\date{}
\begin{document}
\maketitle

\section{Introduction}
E'ma Buyang (ISO 639-3: yzg; Glottocode emab1235) is an endangered Kra-Dai language of the Kra branch. Ethnologue treats it as separate from Langnian Buyang (yln). Search results place it in the Gula River valley, Funing County, Wenshan Prefecture, southeastern Yunnan, China. Estimates differ: about 700 speakers (Ethnologue 2016, via Joshua Project) and about 1,500 for the wider Buyang cluster (1997--2000, per Wikipedia). I found no source placing E'ma Buyang itself in Guangxi. The search found no teaching materials.

\section{Primary Educational Resources}
No primers, textbooks or school programmes for E'ma Buyang turned up in the search.

\section{Linguistic and Reference Materials}
\begin{itemize}
\item Overview of the Buyang language cluster: \url{https://en.wikipedia.org/wiki/Buyang_language}
\item Overview of the Kra languages: \url{https://en.wikipedia.org/wiki/Kra_languages}
\item Overview of the neighbouring Paha language: \url{https://en.wikipedia.org/wiki/Paha_language}
\item A PDF from the Minpaku repository (KH\_025\_2\_004) was returned by the search. I did not confirm its content or its relevance to E'ma Buyang: \url{https://minpaku.repo.nii.ac.jp/record/4093/files/KH_025_2_004.pdf}
\end{itemize}
No grammar or dictionary specific to E'ma Buyang appeared in the search results.

\section{Religious and Community Resources}
\begin{itemize}
\item Joshua Project people-group profile, a mission-oriented source: \url{https://joshuaproject.net/people_groups/19799/CH}
\end{itemize}
I found no scripture translation in E'ma Buyang.

\section{Digital and Audio Resources}
No audio, corpus or software for E'ma Buyang was found.

\section{Curriculum Development Context}
The language is thought to be used by older adults only, and it is not known to be taught in schools. Documentation appears thin. A curriculum would first require documentation work with the community in Gula township, and any school use would take place within the Chinese education system.
\end{document}